\documentclass[aspectratio=169,11pt,xcolor=table]{beamer}
% ------------------------------------------------------------------
% Presentación: A Fast Wait-Free Solution to Read-Reclaim Races
% in Reference Counting (de Wolff, Anderson, Keller, Seletskiy,
% Euro-Par 2024). Compilar con pdfLaTeX.
% Las animaciones son capas de Beamer: avanzan con cada clic
% en modo presentación.
% ------------------------------------------------------------------
\usepackage[utf8]{inputenc}
\usepackage[T1]{fontenc}
\usepackage[spanish,es-noshorthands,es-tabla,es-nodecimaldot]{babel}
\usepackage{lmodern}
\usepackage{booktabs}
\usepackage{pifont}
\usepackage{amsmath,amssymb}
\usepackage{algorithm}
\usepackage[noend]{algpseudocode}
\usepackage{pgfplots}
\pgfplotsset{compat=1.18}
\usetikzlibrary{arrows.meta,positioning,calc,fit,backgrounds,decorations.pathreplacing}

% ---------- Estilo ----------
\definecolor{azul}{HTML}{1F3A5F}
\definecolor{acento}{HTML}{C0392B}
\definecolor{suave}{HTML}{EEF2F7}
\definecolor{gris}{HTML}{6B7280}
\definecolor{cOur}{HTML}{D62728}
\definecolor{cLib}{HTML}{1F77B4}
\definecolor{cWil}{HTML}{2CA02C}
\definecolor{cFol}{HTML}{FF7F0E}
\definecolor{cAnd}{HTML}{9467BD}
\definecolor{amarillo}{HTML}{F6C343}

\usetheme{default}
\usecolortheme{default}
\setbeamercolor{structure}{fg=azul}
\setbeamercolor{frametitle}{fg=white,bg=azul}
\setbeamercolor{title}{fg=azul}
\setbeamercolor{alerted text}{fg=acento}
\setbeamercolor{block title}{fg=white,bg=azul}
\setbeamercolor{block body}{bg=suave}
\setbeamercolor{block title alerted}{fg=white,bg=acento}
\setbeamercolor{block body alerted}{bg=acento!8}
\setbeamerfont{frametitle}{size=\large,series=\bfseries}
\setbeamertemplate{navigation symbols}{}
\setbeamertemplate{itemize items}[circle]
\setbeamertemplate{blocks}[rounded][shadow=false]
\setbeamertemplate{frametitle}{%
  \nointerlineskip
  \begin{beamercolorbox}[wd=\paperwidth,ht=2.4ex,dp=1.2ex,leftskip=1em]{frametitle}
    \usebeamerfont{frametitle}\insertframetitle
  \end{beamercolorbox}}
\setbeamertemplate{footline}{%
  \hfill{\color{gris}\scriptsize\insertframenumber/\inserttotalframenumber}\hspace*{1em}\vspace*{1ex}}

\newcommand{\our}{\textcolor{cOur}{\textbf{Nuestro}}}

% Estilo común para los gráficos de throughput
\pgfplotsset{
  tput/.style={
    width=\linewidth, height=0.62\textheight,
    xmin=-5, xmax=222, ymin=0,
    xtick={1,72,144,216},
    xlabel={Hilos}, ylabel={Mop/s},
    grid=major, grid style={gray!20},
    tick label style={font=\scriptsize}, label style={font=\scriptsize},
    title style={font=\small\bfseries},
    legend style={font=\scriptsize, draw=none, fill=white, fill opacity=0.85, text opacity=1},
    every axis plot/.append style={thick, mark size=1.6pt},
    extra x ticks={144}, extra x tick labels={},
    extra x tick style={grid=major, grid style={dashed, gray!70}},
  },
  sOur/.style={color=cOur, mark=*},
  sLib/.style={color=cLib, mark=square*},
  sWil/.style={color=cWil, mark=x},
  sFol/.style={color=cFol, mark=+},
  sAnd/.style={color=cAnd, mark=star},
}

\title{A Fast Wait-Free Solution to Read-Reclaim Races\\ in Reference Counting}
\subtitle{Conteo de referencias concurrente, \emph{wait-free} y sin reclamación diferida}
\author{Ivo Gabe de Wolff \and Daniel Anderson \and Gabriele K.\ Keller \and Aleksei Seletskiy}
\institute{Euro-Par 2024 \textperiodcentered{} LNCS 14803, pp.\ 103--118\\[2pt]
  {\scriptsize Presentación del artículo}}
\date{}

\begin{document}

% ==================================================================
\begin{frame}[plain]
  \titlepage
\end{frame}

% ==================================================================
\begin{frame}{Contenido}
  \begin{columns}[T]
    \column{0.48\textwidth}
    \begin{enumerate}
      \item<1-> Motivación y problema
      \item<1-> Trabajo previo
      \item<1-> Idea: conteo con pesos
      \item<1-> Algoritmos \textsc{load} / \textsc{store}
    \end{enumerate}
    \column{0.48\textwidth}
    \begin{enumerate}\setcounter{enumi}{4}
      \item<1-> Aumentar el peso sin esperas
      \item<1-> Corrección y complejidad
      \item<1-> Extensiones y límites
      \item<1-> Evaluación experimental
    \end{enumerate}
  \end{columns}
\end{frame}

% ==================================================================
\section{Motivación}
\begin{frame}{Motivación: conteo de referencias}
  \begin{itemize}
    \item<1-> Gestión de memoria segura y simple: cada objeto lleva un \alert{contador de referencias} (\emph{refcount}).
    \item<2-> Muy usado: \texttt{shared\_ptr} en C++, \texttt{Arc} en Rust, implícito en Swift.
    \item<3-> El caso difícil: una \alert{variable mutable compartida} que contiene un puntero con conteo:
      \texttt{atomic<shared\_ptr<T>{}>}.
    \item<4-> Operaciones: \textsc{load} (leer y obtener una referencia propia) y \textsc{store} (reemplazar el puntero).
  \end{itemize}
  \vspace{2mm}
  \onslide<5->{
  \begin{block}{Pregunta del artículo}
    ¿Se puede implementar \texttt{atomic<shared\_ptr>} \alert{wait-free}, \alert{rápido} y \alert{sin sobrecoste de memoria}?
  \end{block}}
\end{frame}

% ==================================================================
\begin{frame}{El problema: carrera \emph{read-reclaim}}
  \centering
  \resizebox{0.97\linewidth}{!}{%
  \begin{tikzpicture}[
      font=\small,
      lane/.style={draw=gray!50, fill=gray!5, rounded corners, minimum height=9mm},
      ev/.style={draw, rounded corners, fill=white, inner sep=3pt, font=\scriptsize, align=center},
      >=Stealth]
    % carriles
    \node[lane, minimum width=13.2cm, anchor=west] (L) at (0,1.6) {};
    \node[lane, minimum width=13.2cm, anchor=west] (W) at (0,0) {};
    \node[anchor=east, font=\small\bfseries] at (-0.1,1.6) {Lector};
    \node[anchor=east, font=\small\bfseries] at (-0.1,0) {Escritor};
    \draw[->, gray] (0,-1.0) -- (13.2,-1.0) node[right, font=\scriptsize] {tiempo};

    \onslide<2->{\node[ev, fill=cLib!10] (a) at (1.6,1.6) {\textbf{1.} lee la variable\\ obtiene ptr a $A$};}
    \onslide<3->{\node[ev] (b) at (4.6,0) {\textbf{2.} \textsc{store}($B$)\\ la variable apunta a $B$};}
    \onslide<4->{\node[ev] (c) at (7.9,0) {\textbf{3.} rc($A$)\,$-1$ = 0\\ \alert{libera $A$}};}
    \onslide<5->{\node[ev, fill=acento!10, draw=acento] (d) at (11.4,1.6) {\textbf{4.} rc($A$)\,$+1$\\ \alert{use-after-free}};}

    \onslide<2->{\draw[densely dotted, gray] (1.6,1.1) -- (1.6,-1.0);}
    \onslide<5->{\draw[densely dotted, gray] (11.4,1.1) -- (11.4,-1.0);
      \draw[decorate, decoration={brace, amplitude=5pt, mirror}, acento, thick]
        (1.6,-1.25) -- node[below=5pt, font=\scriptsize] {ventana peligrosa: leer el puntero e incrementar el contador son \textbf{dos} pasos} (11.4,-1.25);}
  \end{tikzpicture}}
  \vspace{2mm}
  \onslide<6->{\small No existe instrucción que lea la variable \emph{y} modifique el refcount de otro objeto a la vez (DCAS no existe en hardware actual).}
\end{frame}

% ==================================================================
\section{Trabajo previo}
\begin{frame}{Soluciones existentes}
  \centering\small
  \begin{tabular}{@{}lccc@{}}
    \toprule
    \textbf{Enfoque} & \textbf{Wait-free} & \textbf{Sin sobrecoste de memoria} & \textbf{Hardware actual} \\
    \midrule \pause
    Candados (libstdc++)              & \textcolor{acento}{\ding{55}}  & \checkmark & \checkmark \\ \pause
    Lock-free clásico (Valois; Michael \& Scott) & \textcolor{acento}{\ding{55}} & \textcolor{acento}{\ding{55}} & \checkmark \\ \pause
    Detlefs et al.\ (DCAS)            & \textcolor{acento}{\ding{55}}  & \checkmark & \textcolor{acento}{\ding{55}} \\ \pause
    Reclamación diferida (HP, épocas, CDRC) & --- & \textcolor{acento}{\ding{55}} & \checkmark \\ \pause
    Sundell                          & \checkmark & \checkmark & \checkmark\,\footnotesize($\Theta(P)$ por update) \\ \pause
    \emph{Split reference count} (Williams, Folly) & \textcolor{acento}{\ding{55}} & \checkmark & \checkmark \\ \pause
    \alert{\textbf{Este trabajo}} & \checkmark & \checkmark & \checkmark \\
    \bottomrule
  \end{tabular}
  \vspace{3mm}

  \onslide<9->{\small \textsc{store} en $O(1)$; \textsc{load} casi siempre $O(1)$ y $O(P)$ en el peor caso.}
\end{frame}

% ==================================================================
\begin{frame}{Contribuciones}
  \begin{enumerate}
    \item<1-> Estrategia para \texttt{atomic<shared\_ptr>} \alert{wait-free} en \textsc{store}, \textsc{load} y \textsc{swap}, con hardware actual.
    \item<2-> \textsc{load}: normalmente \alert{una sola instrucción atómica}; peor caso $O(P)$.\\
              \textsc{store}: tiempo constante.
    \item<3-> Nueva explicación del \emph{split reference count} como caso particular del \alert{conteo de referencias con pesos}.
    \item<4-> Soporte para compare-and-swap y punteros \emph{aliased} (estos, solo lock-free).
    \item<5-> Implementación en C++ y evaluación frente a cuatro alternativas.
  \end{enumerate}
\end{frame}

% ==================================================================
\section{Idea}
\begin{frame}{Idea clave: conteo de referencias con pesos}
  \begin{columns}[c]
    \column{0.5\textwidth}
    \begin{itemize}
      \item<1-> Bevan; Watson \& Watson (1987).
      \item<2-> Cada referencia tiene un \alert{peso} $\geq 1$.
      \item<3-> Invariante: $\text{rc(obj)} = \sum \text{pesos}$.
      \item<4-> \alert{Copiar} = partir el peso. ¡No toca el refcount!
      \item<5-> Liberar = restar su peso al refcount.
    \end{itemize}
    \column{0.48\textwidth}
    \centering
    \begin{tikzpicture}[font=\small, >=Stealth,
        ref/.style={draw, fill=amarillo, rounded corners=2pt, minimum width=1.5cm, minimum height=6mm},
        obj/.style={draw, fill=amarillo!40, minimum width=2.2cm, minimum height=8mm}]
      \node[obj] (o) at (0,0) {rc = 10};
      \only<1-3>{\node[ref] (x) at (0,2) {$X$: 10}; \draw[->] (x) -- (o);}
      \only<4->{\node[ref] (x) at (-1.2,2) {$X$: 6}; \draw[->] (x) -- (o);
               \node[ref, fill=cFol!60] (y) at (1.2,2) {$Y$: 4}; \draw[->] (y) -- (o);
               \node[font=\scriptsize, gris] at (0,2.75) {6 + 4 = 10 \checkmark};}
      \only<5->{\node[font=\scriptsize, acento, align=center] at (0,-1.1) {tirar $Y$ $\Rightarrow$ rc $-=4$};}
    \end{tikzpicture}
  \end{columns}
\end{frame}

% ==================================================================
\begin{frame}{Pesos en el puntero atómico compartido}
  \begin{columns}[c]
    \column{0.42\textwidth}
    \begin{itemize}
      \item<1-> Al crear el puntero atómico: peso $= C \gg P$ y se suma $C$ al refcount.
      \item<2-> \textsc{load} \alert{resta 1 al peso} y entrega un \texttt{shared\_ptr} de peso 1.
      \item<3-> El refcount del objeto \alert{no cambia}: ya estaba \emph{prepagado}.
    \end{itemize}
    \column{0.56\textwidth}
    \centering
    \begin{tikzpicture}[font=\scriptsize, >=Stealth,
        hdr/.style={draw, fill=amarillo, minimum width=2.6cm, minimum height=5mm},
        hdro/.style={draw, fill=cFol!70, minimum width=2.2cm, minimum height=5mm},
        cell/.style={draw, minimum width=2.6cm, minimum height=5mm}]
      \node[hdr] (ah) at (0,2.2) {Atomic shared pointer};
      \node[cell, below=0pt of ah] (av) {\only<1>{104}\only<2->{\alert{103}}};
      \node[hdr, minimum width=1.9cm] (s1) at (-2.4,1.0) {Shared pointer};
      \node[hdr, minimum width=1.9cm] (s2) at (-2.4,-0.2) {Shared pointer};
      \node[draw, fill=amarillo!40, minimum width=2.2cm, minimum height=5mm] (oh) at (1.4,0.2) {Objeto};
      \node[draw, below=0pt of oh, minimum width=2.2cm, minimum height=5mm] (rc) {Refcount \textbf{106}};
      \draw[->] (av.east) to[bend left=20] (oh.north east);
      \draw[->] (s1.east) -- (oh.west);
      \draw[->] (s2.east) -- (oh.west);
      \onslide<2->{\node[hdro] (n) at (2.9,2.2) {nuevo \texttt{shared\_ptr}};
        \draw[->, cFol, thick] (n.south) -- (oh.north);}
      \node[gris, align=center] at (0,-1.4) {\only<1>{$104 + 1 + 1 = 106$}\only<2->{$103 + 1 + 1 + \alert{1} = 106$ \checkmark}};
    \end{tikzpicture}
  \end{columns}
\end{frame}

% ==================================================================
\begin{frame}{Empaquetar puntero y peso en una palabra}
  \centering
  \begin{tikzpicture}[font=\small, cell/.style={draw, minimum height=7mm}]
    \node[anchor=east] at (-0.2,2) {Puntero};
    \node[cell, minimum width=2.6cm, anchor=west] (z) at (0,2) {$0\ldots0$};
    \node[cell, minimum width=7.2cm, anchor=west, fill=amarillo!60] at (z.east) (ad) {dirección (48 bits)};
    \node[cell, minimum width=1.2cm, anchor=west] at (ad.east) {000};

    \onslide<2->{
    \node[anchor=east] at (-0.2,0.9) {Empaquetado};
    \node[cell, minimum width=2.6cm, anchor=west, fill=cFol!70] (w) at (0,0.9) {peso (16)};
    \node[cell, minimum width=7.2cm, anchor=west, fill=amarillo!60] at (w.east) (ad2) {dirección};
    \node[cell, minimum width=1.2cm, anchor=west] at (ad2.east) {000};}

    \onslide<3->{
    \node[anchor=east] at (-0.2,-0.2) {Con bits de alineación};
    \node[cell, minimum width=3.8cm, anchor=west, fill=cFol!70] (w3) at (0,-0.2) {peso (19)};
    \node[cell, minimum width=7.2cm, anchor=west, fill=amarillo!60] at (w3.east) {dirección};}
  \end{tikzpicture}
  \vspace{4mm}
  \begin{itemize}
    \item<2-> CPUs de 64 bits usan 48 bits de dirección $\Rightarrow$ \alert{16 bits libres}.
    \item<3-> Con los bits de alineación: $C = 2^{19}-1$. En los experimentos: $C = 2^{19}$, $P = 2^{16}$.
    \item<4-> \textsc{load} = un único \alert{fetch-and-add} sobre la palabra: decrementa el peso y devuelve el puntero.
  \end{itemize}
\end{frame}

% ==================================================================
\section{Algoritmos}
\begin{frame}{Algoritmo 1: \textsc{store} y \textsc{load}}
  \begin{columns}[c]
    \column{0.6\textwidth}
    \footnotesize
    \begin{algorithmic}[1]
      \Function{store}{$this$, $new$}
        \State \textcolor<2>{acento}{Incrementar rc($new$) en $C$}
        \State \textcolor<3>{acento}{$(p, w) \gets$ \textproc{unpack}(\textproc{swap}($this$, \textproc{pack}($new$, $C$)))}
        \State \textcolor<4>{acento}{Decrementar rc($p$) en $w$}
      \EndFunction
      \Statex
      \Function{load}{$this$}
        \State \textcolor<5>{acento}{$(p, w) \gets$ \textproc{unpack}(\textproc{fetch\_add}($this$, $-$\textproc{pack}(null, 1)))}
        \If{\textcolor<6>{acento}{$w - 1 < P$}}
          \State \Call{increase\_weight}{$this$, $p$, $w-1$}
        \EndIf
        \State \Return \texttt{shared\_ptr}($p$)
      \EndFunction
    \end{algorithmic}
    \column{0.38\textwidth}
    \small
    \only<1>{Dos operaciones, cada una con \alert{una} instrucción atómica central.}
    \only<2>{Primero se \emph{prepaga} el nuevo objeto con $C$ unidades.}
    \only<3>{Intercambio atómico: se instala $(new, C)$ y se recupera el par anterior.}
    \only<4>{Solo se restan los $w$ tickets \emph{no usados}; los lectores en curso siguen contando.}
    \only<5>{Leer y «apuntarse» es \alert{un solo paso}: no hay ventana para la carrera.}
    \only<6>{Si el peso cae por debajo de $P$, este hilo debe recargarlo antes de continuar.}
  \end{columns}
\end{frame}

% ==================================================================
\begin{frame}{Aumentar el peso: requisitos}
  \begin{block}{Problema}
    Recargar el peso no puede ser un fetch-and-add: hay que comprobar que el puntero no cambió ni fue ya recargado $\Rightarrow$ bucle de compare-and-swap.
  \end{block}
  \vspace{2mm}
  \onslide<2->{
  \begin{block}{R-1}
    $C \geq 2P$: el peso máximo es al menos el doble del número de hilos.
  \end{block}}
  \onslide<3->{
  \begin{block}{R-2}
    El peso solo cambia de dos formas: \alert{baja de uno en uno} o \alert{salta por encima de $P$}.
  \end{block}}
  \vspace{2mm}
  \onslide<4->{\small\textbf{Consecuencia:} si un CAS falla porque el peso \emph{subió}, el peso estuvo por encima de $P$ $\Rightarrow$ otro hilo ya recargó $\Rightarrow$ \alert{salida temprana}.}
\end{frame}

% ==================================================================
\begin{frame}{Algoritmo 2: \textsc{increase\_weight} (wait-free)}
  \begin{columns}[c]
    \column{0.62\textwidth}
    \footnotesize
    \begin{algorithmic}[1]
      \Function{increase\_weight}{$value$, $ptr$, $weight$}
        \State \textcolor<2>{acento}{Incrementar rc($ptr$) en $C - weight$}
        \State $cur \gets weight$
        \Loop
          \State \textcolor<3>{acento}{$cas \gets$ \Call{CAS}{$value$, $(ptr, cur)$, $(ptr, cur + C - weight)$}}
          \If{$cas.success$} \Return \EndIf
          \State $(p', w') \gets$ \Call{unpack}{$cas.original$}
          \If{\textcolor<4>{acento}{$p' \neq ptr$ \textbf{or} $w' > cur$}} \textbf{break} \EndIf
          \State \textcolor<5>{acento}{$cur \gets w'$}
        \EndLoop
        \State Decrementar rc($ptr$) en $C - weight$
      \EndFunction
    \end{algorithmic}
    \column{0.36\textwidth}
    \small
    \only<1>{Bucle CAS con \alert{tres} salidas posibles.}
    \only<2>{\alert{Reservar antes de publicar}: nunca hay pesos en el puntero sin respaldo en el refcount.}
    \only<3>{Instala el nuevo peso solo si nada cambió.}
    \only<4>{\textbf{Salida temprana}: otro objeto, o el peso subió (R-2) $\Rightarrow$ ya recargado.}
    \only<5>{Si el peso bajó, reintenta. Los valores esperados decrecen estrictamente: $\leq P$ iteraciones.}
  \end{columns}
\end{frame}

% ==================================================================
\begin{frame}{El peso a lo largo del tiempo}
  \centering
  \begin{tikzpicture}[>=Stealth, font=\small, x=1.15cm, y=1cm]
    \fill[amarillo!25] (0,0) rectangle (10,2);
    \draw[->] (0,0) -- (10.3,0) node[right] {tiempo};
    \draw[->] (0,0) -- (0,4) node[above] {peso};
    \draw[dashed, gris] (0,2) -- (10,2) node[right] {$P$};
    \node[left] at (0,2) {$P$};
    \node[left] at (0,0) {0};
    % curva (por tramos)
    \onslide<1->{\draw[very thick, cFol] (0,1.2) -- (1,0.9) -- (2,0.6) -- (2.6,0.4);}
    \onslide<2->{\draw[very thick, cFol] (2.6,0.4) -- (2.6,3.1) -- (3.6,2.7) -- (4.6,2.2) -- (5.4,1.6) -- (6.0,1.4);}
    \onslide<4->{\draw[very thick, cFol] (6.0,1.4) -- (6.4,1.3) -- (6.4,3.6) -- (7.6,3.3) -- (9,3.0) -- (10,2.8);}
    % marcas
    \onslide<1->{\draw[dashed] (2.2,0) -- (2.2,0.55); \node[below] at (2.2,0) {1};}
    \onslide<2->{\draw[dashed] (3.6,0) -- (3.6,2.7); \node[below] at (3.6,0) {2};}
    \onslide<3->{\draw[dashed] (5.7,0) -- (5.7,1.5); \node[below] at (5.7,0) {3};}
    \onslide<4->{\draw[dashed] (7.6,0) -- (7.6,3.3); \node[below] at (7.6,0) {4};}
  \end{tikzpicture}
  \vspace{3mm}

  \small
  \only<1>{\textbf{1.} Un hilo hace \textsc{load}, deja el peso bajo $P$ y llama a \textsc{increase\_weight}.}
  \only<2>{\textbf{2.} Otro hilo recarga primero: el peso supera $P$.}
  \only<3>{\textbf{3.} Nuestro hilo ve un CAS fallido con peso \alert{mayor} al esperado (aunque bajo $P$) $\Rightarrow$ sale: ya hubo recarga.}
  \only<4>{\textbf{4.} Los hilos que bajaron el peso después de 2 son responsables de la siguiente recarga.}
\end{frame}

% ==================================================================
\section{Corrección}
\begin{frame}{Corrección: el peso nunca es negativo}
  \begin{itemize}
    \item<1-> \textbf{Ofensa}: un hilo decrementa el peso cuando este es $\leq P$.
    \item<2-> $O$: multiconjunto de hilos ofensores; se vacía cuando el peso supera $P$.
  \end{itemize}
  \vspace{1mm}
  \onslide<3->{
  \begin{block}{Invariante 1: cada hilo está en $O$ como mucho una vez}
    Tras ofender, el hilo debe salir de \textsc{increase\_weight}; las tres salidas implican que el peso estuvo $> P$ $\Rightarrow$ $O$ se vació.
  \end{block}}
  \onslide<4->{
  \begin{block}{Invariante 2: peso $\geq P - |O|$}
    Por R-2, desde el último vaciado el peso solo baja de uno en uno, y cada bajada bajo $P$ añade un hilo a $O$.
  \end{block}}
  \onslide<5->{
  \begin{alertblock}{Conclusión}
    $|O| \leq P \;\Rightarrow\;$ peso $\geq P - P = 0$. \quad Con 1 o con $2^{16}$ hilos.
  \end{alertblock}}
\end{frame}

% ==================================================================
\begin{frame}{Complejidad}
  \centering
  \begin{tabular}{@{}lll@{}}
    \toprule
    \textbf{Operación} & \textbf{Progreso} & \textbf{Coste} \\
    \midrule
    \textsc{load}   & wait-free & casi siempre $O(1)$; peor caso $O(P)$ \\
    \textsc{store}  & wait-free & $O(1)$ \\
    \textsc{swap}   & wait-free & $O(1)$ \\
    compare-and-swap & lock-free & --- \\
    punteros \emph{aliased} & wait-free (rápido) / lock-free (fallback) & --- \\
    \bottomrule
  \end{tabular}
  \vspace{4mm}
  \begin{itemize}
    \item<2-> En la práctica, el peor caso depende de los hilos que \alert{realmente} usan esa variable, no de $P$.
    \item<3-> Con $C=2^{19}$ y $P=2^{16}$: $\approx 4{,}6\times10^5$ \textsc{load}s entre recargas.
  \end{itemize}
\end{frame}

% ==================================================================
\section{Extensiones}
\begin{frame}{Extensiones y limitaciones}
  \begin{columns}[T]
    \column{0.48\textwidth}
    \textbf{Extensiones}
    \begin{itemize}
      \item<1-> \textbf{null}: el peso puede desbordar sin alterar el puntero (bits altos).
      \item<2-> \textbf{Punteros débiles}: ortogonales; mismo esquema.
      \item<3-> \textbf{Compare-and-swap}: lock-free (falla si cambia el peso).
      \item<4-> \textbf{Aliased}: palabra de 128 bits; ruta rápida optimista + fallback lock-free.
    \end{itemize}
    \column{0.48\textwidth}
    \onslide<5->{\textbf{Limitaciones}
    \begin{itemize}
      \item<5-> Límite duro de $P$ hilos.
      \item<6-> Requiere bits libres en los punteros (afecta tagging, LA57, ARM TBI/PAC).
      \item<7-> Cada \textsc{load} escribe en la variable: contención en lecturas.
      \item<8-> Sin \emph{snapshot pointers}.
      \item<9-> Refcount de 64 bits recomendado.
    \end{itemize}}
  \end{columns}
\end{frame}

% ==================================================================
\section{Evaluación}
\begin{frame}{Configuración experimental}
  \begin{columns}[c]
    \column{0.52\textwidth}
    \begin{itemize}
      \item<1-> \textbf{Throughput}: 4$\times$ Intel Xeon E7-8867, 72 núcleos / \alert{144 hilos SMT}, 1\,TB.
      \item<2-> \textbf{Latencia}: Intel Xeon w7-3465X, 28 núcleos / \alert{56 hilos}, 128\,GB.
      \item<3-> Pila concurrente: $n$ pilas, $u$\,\% de actualizaciones.
      \item<4-> $C = 2^{19}$, $P = 2^{16}$.
    \end{itemize}
    \column{0.44\textwidth}
    \onslide<5->{
    \textbf{Implementaciones comparadas}\\[1mm]
    \small
    \textcolor{cOur}{$\bullet$} \textbf{Nuestro} (wait-free)\\
    \textcolor{cLib}{$\bullet$} \textbf{libstdc++} (candados)\\
    \textcolor{cWil}{$\bullet$} \textbf{Williams} (split rc, lock-free)\\
    \textcolor{cFol}{$\bullet$} \textbf{Folly} (split rc, lock-free)\\
    \textcolor{cAnd}{$\bullet$} \textbf{Anderson} (hazard pointers;\\ \hspace*{0.8em}pequeño sobrecoste de memoria)}
  \end{columns}
\end{frame}

% ==================================================================
\begin{frame}{Throughput: baja contención ($n=10$, $u=1\%$)}
  \centering
  \begin{tikzpicture}
    \begin{axis}[tput, width=0.82\linewidth, height=0.72\textheight, ymax=52,
        legend pos=outer north east]
      \only<1->{\addplot[sLib] coordinates {(1,6.7)(18,4.3)(36,2.5)(54,1.8)(72,1.4)(108,1.2)(144,1.0)(180,0.8)(216,0.6)}; \addlegendentry{libstdc++}}
      \only<1->{\addplot[sWil] coordinates {(1,4)(18,14.5)(36,16)(54,16)(72,13.3)(108,9.5)(144,6.5)(180,6.6)(216,7)}; \addlegendentry{Williams}}
      \only<1->{\addplot[sFol] coordinates {(1,2)(18,10)(36,5.5)(54,4.3)(72,3.8)(108,3.3)(144,3.0)(180,2.8)(216,2.8)}; \addlegendentry{Folly}}
      \only<2->{\addplot[sAnd] coordinates {(1,5.5)(18,25.5)(36,31)(54,34.5)(72,36.8)(108,32.3)(144,30.4)(180,30.4)(216,30.4)}; \addlegendentry{Anderson}}
      \only<3->{\addplot[sOur, very thick] coordinates {(1,9.3)(18,30)(36,38.5)(54,44)(72,47)(108,46)(144,46.5)(180,46)(216,46.5)}; \addlegendentry{Nuestro}}
    \end{axis}
  \end{tikzpicture}

  \small
  \only<1>{Los métodos sin sobrecoste de memoria colapsan al crecer los hilos.}
  \only<2>{Anderson (reclamación diferida) escala mejor\ldots}
  \only<3>{\ldots pero el método propuesto llega a $\approx$\,\alert{47 Mop/s}, $\sim$1{,}5$\times$ Anderson y $\sim$7$\times$ Williams.}
\end{frame}

% ==================================================================
\begin{frame}{Throughput: alta contención vs.\ muchas pilas}
  \begin{columns}[T]
    \column{0.5\textwidth}
    \centering
    \begin{tikzpicture}
      \begin{axis}[tput, title={$n=1$, $u=1\%$}, ymax=11.5, legend style={at={(0.5,1.18)}, anchor=south, legend columns=3}]
        \addplot[sLib] coordinates {(1,7.0)(18,0.5)(36,0.3)(54,0.2)(72,0.15)(108,0.1)(144,0.1)(180,0.1)(216,0.1)};
        \addplot[sWil] coordinates {(1,3.3)(18,2.2)(36,1.9)(54,1.6)(72,1.4)(108,1.05)(144,0.8)(180,0.8)(216,0.8)};
        \addplot[sFol] coordinates {(1,1.8)(18,0.7)(36,0.55)(54,0.45)(72,0.6)(108,0.45)(144,0.35)(180,0.35)(216,0.3)};
        \addplot[sAnd] coordinates {(1,5.4)(18,4.7)(36,6.1)(54,7.5)(72,8.3)(108,9.7)(144,10.0)(180,10.4)(216,10.3)};
        \addplot[sOur, very thick] coordinates {(1,9.1)(18,6.6)(36,6.9)(54,6.3)(72,6.3)(108,6.1)(144,5.85)(180,6.0)(216,5.8)};
      \end{axis}
    \end{tikzpicture}
    \column{0.5\textwidth}
    \centering
    \onslide<2->{
    \begin{tikzpicture}
      \begin{axis}[tput, title={$n=10^6$, $u=50\%$}, ymax=140]
        \addplot[sLib] coordinates {(1,1)(18,18.5)(36,38)(54,55.5)(72,73.5)(108,94.5)(144,112.5)(180,108.5)(216,103)};
        \addplot[sWil] coordinates {(1,0.8)(18,15)(36,28)(54,40)(72,53)(108,65)(144,76)(180,75.5)(216,76)};
        \addplot[sFol] coordinates {(1,0.5)(18,10)(36,18)(54,26)(72,35)(108,43)(144,51.5)(180,50.5)(216,51.5)};
        \addplot[sAnd] coordinates {(1,1)(18,16)(36,29)(54,41.5)(72,53.5)(108,65.5)(144,76.5)(180,75)(216,72)};
        \addplot[sOur, very thick] coordinates {(1,1)(18,19)(36,38.5)(54,57)(72,77)(108,104)(144,129)(180,128.5)(216,129)};
      \end{axis}
    \end{tikzpicture}}
  \end{columns}
  \small
  \only<1>{\textcolor{cOur}{\textbf{Nuestro}} \textcolor{gris}{(rojo)} es el mejor sin sobrecoste de memoria, pero con \alert{una sola pila} Anderson gana: todas las lecturas escriben en la misma palabra.}
  \only<2>{Con \alert{$10^6$ pilas} (poca contención) gana \textcolor{cOur}{\textbf{Nuestro}}: $\approx$\,129 Mop/s frente a 113 (libstdc++) y 77 (Anderson).}
\end{frame}

% ==================================================================
\begin{frame}{Latencia (mediana)}
  \begin{columns}[T]
    \column{0.5\textwidth}
    \centering
    \begin{tikzpicture}
      \begin{axis}[
          ybar, bar width=4.5pt, width=\linewidth, height=0.62\textheight,
          title={Sin contención ($p=1$), ns}, title style={font=\small\bfseries},
          symbolic x coords={Load,Store,Store+reclaim}, xtick=data,
          ymin=0, ymax=55, enlarge x limits=0.25,
          tick label style={font=\scriptsize},
          legend style={font=\scriptsize, at={(0.5,1.2)}, anchor=south, legend columns=4, draw=none},
          ymajorgrids, grid style={gray!20}]
        \addplot[fill=cOur, draw=cOur]  coordinates {(Load,23)(Store,33)(Store+reclaim,44)};
        \addplot[fill=cLib, draw=cLib]  coordinates {(Load,29)(Store,36)(Store+reclaim,41)};
        \addplot[fill=cFol, draw=cFol]  coordinates {(Load,24)(Store,37)(Store+reclaim,49)};
        \addplot[fill=cAnd, draw=cAnd]  coordinates {(Load,28)(Store,28)(Store+reclaim,44)};
        \legend{Nuestro, libstdc++, Folly, Anderson}
      \end{axis}
    \end{tikzpicture}
    \column{0.5\textwidth}
    \centering
    \onslide<2->{
    \begin{tikzpicture}
      \begin{axis}[
          ybar, bar width=4.5pt, width=\linewidth, height=0.62\textheight,
          title={Alta contención ($p=56$), $\mu$s, escala log}, title style={font=\small\bfseries},
          symbolic x coords={Load,Store,Store+reclaim}, xtick=data,
          ymode=log, ymin=0.5, ymax=300, log origin=infty, enlarge x limits=0.25,
          tick label style={font=\scriptsize},
          ymajorgrids, grid style={gray!20}]
        \addplot[fill=cOur, draw=cOur]  coordinates {(Load,1.14)(Store,1.47)(Store+reclaim,1.57)};
        \addplot[fill=cLib, draw=cLib]  coordinates {(Load,99.7)(Store,111)(Store+reclaim,59.3)};
        \addplot[fill=cFol, draw=cFol]  coordinates {(Load,11.6)(Store,2.62)(Store+reclaim,5.06)};
        \addplot[fill=cAnd, draw=cAnd]  coordinates {(Load,2.58)(Store,2.84)(Store+reclaim,1.25)};
      \end{axis}
    \end{tikzpicture}}
  \end{columns}
  \small
  \only<1>{Sin contención todos son similares; \textsc{load} más rápido con \textcolor{cOur}{\textbf{Nuestro}} (23 ns).}
  \only<2>{Con contención, \textsc{load} es \alert{$\sim$2{,}3$\times$} más rápido que Anderson y \alert{$\sim$87$\times$} que libstdc++. Store+reclaim: Anderson ligeramente mejor.}
\end{frame}

% ==================================================================
\begin{frame}{Relación con el \emph{split reference count}}
  \begin{columns}[c]
    \column{0.55\textwidth}
    \begin{itemize}
      \item<1-> Split rc: un \alert{contador local} en el puntero que empieza en 0 y \emph{sube} con cada lectura.
      \item<2-> Aquí: un \alert{peso} que empieza en $C$ y \emph{baja}.
      \item<3-> Equivalencia: $\text{peso} = C - \text{contador local}$.
      \item<4-> Diferencia clave: \alert{salida temprana} en \textsc{increase\_weight} $\Rightarrow$ wait-free.
      \item<5-> Folly: refcount de 32 bits, $C=8192$, CAS en \textsc{load}; puede volverse bloqueante con $>C$ hilos.
    \end{itemize}
    \column{0.42\textwidth}
    \centering
    \begin{tikzpicture}[font=\scriptsize]
      \begin{axis}[width=\linewidth, height=0.55\textheight, xmin=0, xmax=8, ymin=0, ymax=8.5,
          xlabel={lecturas}, xtick={0,2,...,8}, ytick={0,4,8}, yticklabels={0,,$C$},
          legend style={font=\scriptsize, draw=none, at={(0.5,1.02)}, anchor=south, legend columns=2},
          tick label style={font=\scriptsize}]
        \only<1->{\addplot[thick, cLib, mark=*, mark size=1pt] coordinates {(0,0)(8,8)}; \addlegendentry{contador local}}
        \only<2->{\addplot[thick, cOur, mark=*, mark size=1pt] coordinates {(0,8)(8,0)}; \addlegendentry{peso}}
      \end{axis}
    \end{tikzpicture}
  \end{columns}
\end{frame}

% ==================================================================
\begin{frame}{Conclusiones}
  \begin{itemize}
    \item<1-> \texttt{atomic<shared\_ptr>} puede ser \alert{wait-free sin reclamación diferida}.
    \item<2-> La carrera read-reclaim se elimina asignando un \alert{peso} al puntero atómico: leer y adquirir la referencia es una sola instrucción.
    \item<3-> R-1 ($C\geq2P$) + R-2 + salida temprana $\Rightarrow$ recarga acotada por $P$.
    \item<4-> En la práctica: \alert{mayor throughput} que las alternativas sin sobrecoste y \alert{latencia de \textsc{load} mínima}.
    \item<5-> Límites: CAS solo lock-free, contención en lecturas de una misma variable, dependencia de bits libres.
  \end{itemize}
  \vspace{3mm}
  \onslide<6->{\small\textcolor{gris}{Código y artefacto: Zenodo (doi:10.5281/zenodo.11775182) y github.com/ivogabe/wait-free-atomic-shared-pointer}}
\end{frame}

% ==================================================================
\begin{frame}[plain]
  \centering
  \vspace*{20mm}
  {\Huge\color{azul} ¡Gracias!}\\[6mm]
  {\large Preguntas}
\end{frame}

\end{document}
